\bookchapter{Bab 0: Hai, Sebelum Mula}

Sebelum meneroka secara rasmi, mari kita bercakap tentang peraturan asas untuk memastikan anda faham apa yang akan dibincangkan kemudian. Jangan tergesa-gesa, mengasah pisau tidak melambatkan kerja menebang kayu.

\booksection{I. Apa yang Buku Ini Bantu Anda Lakukan? (Dan Apa yang Tidak)}

Ia membantu anda memahami secara sistematik cara menggunakan telefon pintar, memahami sebab di sebalik fungsi, dan sama ada anda penyokong Android atau Apple, kebanyakan idea adalah universal. Ia akan membimbing anda berlatih, menaikkan taraf anda dari “tahu guna” ke “faham guna”.

Ia bukan manual untuk model telefon tertentu, bukan berita teknologi mengejar produk baharu, dan bukan panduan mengajar anda ketagih telefon. Sejujurnya, jika anda menjangkakan “cara menonton video dengan lebih puas”, anda mungkin kecewa — kami lebih mengambil berat tentang cara menjadikan telefon benar-benar berkhidmat kepada anda, bukan sebaliknya.

\booksection{II. Penjelasan “Kod Rahsia” Kami}

Untuk menjelaskan dengan lebih jelas, buku ini menggunakan beberapa penanda kecil, penanda ini seperti kod rahsia antara kita:

\begin{itemize}
  \item \textbf{[Nama Butang]}: Mewakili butang atau nama App tertentu pada skrin, contohnya [Tetapan], [OK].
  \item \textbf{>}: Mewakili urutan operasi, contohnya "Ketik [Tetapan] > [Bunyi]" bermaksud ketik Tetapan dahulu, kemudian ketik Bunyi.
  \item \textbf{Tebal}: Sangat penting! Boleh jadi konsep utama atau amaran keselamatan, sila ambil perhatian.
  \item \emph{Italik}: Beberapa idea atau metafora tambahan untuk membantu anda memahami dengan lebih baik.
\end{itemize}

\booksection{III. Android dan iOS, Kita Bincang Kedua-duanya}

Anda tahu apa, sesetengah perkara dipanggil dengan nama berbeza tetapi fungsinya sama. Contohnya: Aplikasi (Android) dan App (iOS), Tetapan (kedua-duanya dipanggil Tetapan), Pengurusan Fail (Android) dan Fail (iOS). Asalkan anda tahu ia perkara yang sama, jangan pening dengan nama.

Selain itu, dalam kem Android kini ada sistem HarmonyOS Huawei, yang mempunyai ciri tersendiri, akan dibincangkan khusus kemudian. Ringkasnya, tidak kira telefon apa yang anda gunakan, buku ini boleh membantu anda.

\booksection{IV. Kapasiti Telefon: Mengapa Beli 128G tetapi Kurang daripada 110G Tersedia?}

Semua orang tahu unit kapasiti: 1TB = 1024GB, 1GB = 1024MB, 1MB = 1024KB... tetapi ada perangkap di sini: pengeluar cakera keras biasanya mengira berdasarkan asas 1000, manakala sistem mengira berdasarkan asas 1024, jadi cakera keras “128G” yang anda beli akan “menyusut” dari segi kapasiti sebenar.

Tapi sebab utama kapasiti telefon anda menyusut ada tiga:

\begin{enumerate}
  \item \textbf{Sistem itu sendiri mengambil ruang}: Sistem pengendalian seperti dinding struktur rumah, mesti wujud. Sistem kini mudah mengambil 10-20GB, ini sangat normal.
  \item \textbf{Aplikasi prat dipasang pengeluar}: App yang disertakan dengan telefon baharu juga mengambil ruang, sesetengah boleh dibuang, sesetengah tidak.
  \item \textbf{Kehilangan pembahagian dan pemformatan}: Atasan teknikal, seperti pelapis beg sekolah juga mengambil isipadu.
\end{enumerate}

\begin{tcbinfo}{Saranan Membeli Telefon}
Telefon 128GB kini pada asasnya tidak mencukupi, 256GB sebagai permulaan lebih selamat. Jika anda suka mengambil gambar, bermain permainan, cadangan terus ke 512GB atau lebih besar. Ingat, ruang storan ini, lebih baik berlebihan daripada tidak mencukupi.

\end{tcbinfo}

\begin{tcbpractice}{Mari Berlatih: Lihat "Kad Pengenalan" Telefon Anda}
Buka [Tetapan], cari Perihal Telefon (atau pilihan serupa), catat: model telefon, versi sistem, jumlah ruang storan dan ruang yang tinggal. Ini adalah langkah pertama untuk mengenalinya. Semak juga kesihatan bateri, supaya anda ada gambaran.

\end{tcbpractice}

\booksection{V. Telefon Anda Adalah “Perpustakaan Aset Digital”}

Sebelum memulakan pengetahuan khusus, mari kita bina persefahaman teras:\textbf{telefon bukan lagi sekadar alat komunikasi, ia adalah perpustakaan aset digital paling penting kami}. Kenalan, rekod sembang, gambar keluarga, dokumen kerja, App bank, kod pembayaran, data kesihatan, dokumen elektronik... data ini digabungkan, nilainya sering melebihi telefon itu sendiri.

\begin{tcbinfo}{Mengapa Ini Penting?}
Jika telefon hilang atau rosak, anda kehilangan bukan sahaja perkakasan bernilai beberapa ribu yuan,\textbf{tetapi lebih besar kemungkinan data tidak boleh diperbaharui yang terkumpul selama bertahun-tahun}— gambar pertumbuhan anak, suara saudara yang telah meninggal, kontrak penting yang tidak disandarkan. Tabiat penyusunan dan perlindungan saintifik boleh melindungi anda dari tiga dimensi:

\begin{itemize}
  \item \textbf{Meningkatkan kecekapan operasi}: Klasifikasi fail yang standard membolehkan sistem mencari lebih cepat, telefon tidak tergantung jangka panjang.
  \item \textbf{Memanjangkan jangka hayat}: Pengurusan storan yang munasabah mengurangkan beban baca-tulis cip, bateri dan storan lebih tahan lama.
  \item \textbf{Menjamin keselamatan data}: Pengurusan berlapis + sandaran berkala + penumpuan kebenaran, menjauhkan "bencana data" daripada anda.
\end{itemize}

\end{tcbinfo}

\booksubsection{Falsafah Teras Kami: “Pencegahan Utama, Gabungan Pencegahan dan Rawatan”}

Seperti menjaga kesihatan tidak boleh menunggu sakit baru jumpa doktor,\textbf{pengurusan telefon juga mesti bermula dari “tabiat harian”}. Buku ini menekankan berulang kali bukan kemahiran lanjutan tertentu, tetapi prinsip-prinsip ringkas ini:

\begin{itemize}
  \item \textbf{Pencegahan utama}: Pasang App hanya dari kedai rasmi, nyahpeka gambar sebelum dihantar, jangan ulang kata laluan — "tindakan kecil" ini kosnya sangat rendah, tetapi boleh menangkis kebanyakan risiko.
  \item \textbf{Pemeriksaan berkala}: Susun folder muat turun sekali sebulan, audit kebenaran sekali setiap suku tahun, sandaran penuh sekali setiap enam bulan — rentetan tetap 100 kali lebih berkesan daripada memadam kebakaran sementara.
  \item \textbf{Berikan ruang untuk berundur}: Aktifkan penyegerakan awan, ingat kata laluan akaun, tulis "senarai warisan digital" — pastikan apabila kemalangan berlaku, anda dan keluarga boleh menghadapi dengan tenang.
\end{itemize}
